\section{Conditional arbitrary-companion transfer from intrinsic harmonic locality}
\label{sec:r24-harmonic-locality-arbitrary-companions}

\paragraph{Status.}
This is a \texttt{CONDITIONAL} deduction with zero new principal
theorem count. Its explicit antecedent is the still open intrinsic
harmonic-locality assertion IL20 for genuine original full polars.
The scalar recognition HL3 and the HC equivalence identify that
antecedent with the corresponding original local-unit drift gate
NC12/HE7 in their printed setting. No proof of that antecedent is
supplied here. The purpose is to replace the global-comparator
shortcut disproved by BC with an exact, sufficient local-measure
interface, while retaining every fixed-domain quantifier.

All fields, norms, spaces and ordinary opens are original. The
spaces \(\mathcal D\), the closed carrier \(F\), compact stage \(L\),
formal symmetric powers, relation generators and seminorms are
exactly KC1--KC5. In particular a full polar is a polar of that
entire seminorm, not a real conormal current selected independently.

\subsection{The precise assumption}

Choose an original ordinary \(U\) in a good affinoid interior.
For every closed carrier \(F\subset U\), every compact
\(L\Subset U\), every degree \(m\ge2\), every genuine original
unit full polar \(\Lambda\) of \(q^{U,m}_{F,L}\), and arbitrary
fixed \(u_3,\ldots,u_m\in\mathcal D(U)\), IR and IL give a finite
signed measure \(\nu_{u,u_3,\ldots,u_m}\) supported on
\(C=F\cap L\), characterized by
\[
 \Lambda([k][u][u_3]\cdots[u_m])
       =\int_C k\,d\nu_{u,u_3,\ldots,u_m}
       \qquad(k,u\in\mathcal D(U)).                         \tag{LC24.1}
\]
We assume, on this same \(U\), the following assertion:
\[
 \begin{gathered}
 u\in\mathcal D(U),\quad W\subset U\text{ ordinary open},\\
 u|_{F\cap W}\text{ has original continuous harmonic trace
 witnesses locally at every point}\\
 \Longrightarrow\quad
 \nu_{u,u_3,\ldots,u_m}|_W=0 .
 \end{gathered}                                             \tag{LC24.2}
\]
All choices in the preceding paragraph are quantified in this
assumption. It concerns restriction of a measure already defined
from \(\Lambda\), not restriction of \(\Lambda\) to a new
localized relation quotient.

\subsection{Strict geometric defect gives local harmonicity without global witnesses}

Let \(\pi:Y\to U_*\) be an original finite \'{e}tale surjection
of degree \(d\ge1\), let \(x\in U_*\), and fix a global analytic
unit \(g\) on \(Y\). GP prepares one ordinary \(U\ni x\),
a power frame \(e\) of the finite algebra, and a finite analytic
gallery from its flag to the flag of \(ge\).
Every relevant prefix and adjacent-exchange coordinate minor is
a unit on the complete actual ordering cover \(q:Z_U\to U\).
This preparation depends on \(\pi,g,x\), not on a later weight
\(w\in\mathcal D(Y_U)\), degree, carrier, compact stage, polar,
or downstream tail. The additional finite coefficient selection
in CW is not used here.

For an adjacent exchange use the notation
\[
 a=W\wedge u,\quad b=W\wedge v,\quad
 \alpha=\log\|a\|_w,\quad \beta=\log\|b\|_w,\quad
 c=\log\|W\|_w,\quad d_0=\log\|W\wedge u\wedge v\|_w,
                                                               \tag{LC24.3}
\]
where the norms on the actual ordering cover are the diagonal
sheet norms with weights pulled back from \(w\). Their prefix
formulas are finite maxima of sums of those weights and logarithms
of the protected analytic minors. FS4--FS5 therefore descend all
four functions as members of \(\mathcal D(U)\), for arbitrary
\(w\in\mathcal D(Y_U)\).
Set
\[
                         H=\alpha-\beta,\qquad
                         K_0=\alpha+\beta-c-d_0\ge0.        \tag{LC24.4}
\]
The last inequality is the exact exterior-norm inequality proved
in FG/CW; it is pointwise on all original analytic points.

\begin{lemma}[Local harmonicity of the strict-defect function]
For every \(w\in\mathcal D(Y_U)\), the descended scalar \(H\)
is original current-harmonic locally on the ordinary open
\(\{K_0>0\}\subset U\). No global harmonic comparator on \(U\)
is asserted.
\end{lemma}

\begin{proof}
Fix an original point \(z\) of the ordering cover above
\(y\in\{K_0>0\}\). The strict exterior inequality in CW18--CW20
applies with the continuous positive weights of this \(w\).
Choose a coordinate \(I\) maximizing the weighted norm of \(a\).
The strict determinant inequality makes it a maximizer for \(b\)
as well, and gives
\[
          \|a/a_I-b/b_I\|_w<q_I
\]
at \(z\), where \(q_I\) is the weight of that coordinate.
Both normalized vectors have \(I\)-coordinate \(1\).
The strict inequality persists on an ordinary neighborhood of
\(z\), since the weights are continuous and \(a_I,b_I\) are
protected analytic units. The ultrametric inequality then gives
equality of the two normalized vector norms throughout that
neighborhood. Consequently
\[
                          q^*H=\log|a_I/b_I|               \tag{LC24.5}
\]
there. The exterior-norm identity used to pass from the prefix
defect to this inequality is the same finite Gaussian-elimination
identity in CW18; it uses no PL presentation of \(w\).

There are finitely many actual points over \(y\). Apply (LC24.5)
near each of them. Finite properness allows one ordinary
neighborhood \(V\ni y\) whose complete inverse image lies in
the union of these neighborhoods. Thus \(q^*H\) and
\(-q^*H\) are classically psh on all of \(q^{-1}(V)\).
The finite-surjective scalar descent theorem used in FS, applied
separately to the two signs, makes \(H|_V\) classically
pluriharmonic. HL3 identifies it with an original continuous
current-harmonic function, and indeed with a local original
unit-logarithm sum. This is a scalar descent on the complete
finite inverse image. No arbitrary current has been pulled
through the finite map, and no full polar has been localized.
\end{proof}

\subsection{Conditional exchange with every downstream tail retained}

Assume (LC24.2) on the prepared \(U\). For any full polar and
any fixed additional \(\mathcal D(U)\)-tails, the preceding
lemma and (LC24.2) show that the measure with distinguished tail
\(H\) is zero on \(\{K_0>0\}\).
Its remaining support is contained in \(\{K_0=0\}\).
Equation (LC24.1) therefore gives
\[
                 \Lambda([K_0][H]R)=0                    \tag{LC24.6}
\]
for every formal monomial \(R\) in arbitrary downstream
\(\mathcal D(U)\)-factors. The total degree is at least \(2\);
the measure is finite and \(K_0\) is continuous, so this
integration on its support is legitimate.

For \(b\ge2\), put \(A=\alpha-c\), \(B=\beta-c\) and
\(T=d_0-c\). The real polynomial
\[
 A^b+(T-A)^b-B^b-(T-B)^b
\]
is divisible by \((A-B)(A+B-T)\). Indeed it vanishes at
\(A=B\) and at \(A=T-B\); these are relatively prime linear
polynomials in \(\mathbf R[A,B,T]\).
Their product is \(H K_0\), and the quotient is homogeneous of
degree \(b-2\). This is an exact formal polynomial identity
in the degree-one space \(\mathcal D(U)\).
Expanding the quotient into finitely many monomials, multiplying
by any extra monomial \(R\), and applying (LC24.6) proves that
every full polar annihilates the adjacent power-sum difference
times \(R\). CV2 yields
\[
 q^{U,b+\deg R}_{F,L}
 \left(
   \bigl([\alpha-c]^b+[d_0-\alpha]^b
        -[\beta-c]^b-[d_0-\beta]^b\bigr)R
 \right)=0.                                                 \tag{LC24.7}
\]
For \(b=1\) the displayed difference is literally zero.
At joins between two analytic adapted frames for the same flag,
the prefix logs change by logarithms of global base unit
determinants. Expanding powers then gives the exact harmonic
factor certificates, including all factors in \(R\).
Summing (LC24.7) along the fixed gallery gives UT5, and hence
UT7, for every \(w\in\mathcal D(Y_U)\), without a fixed
finite-unit PL restriction. In degree \(d=1\), this conclusion
already follows from the global unit factor and needs no exchange.

\subsection{The original KC limit still has a uniform budget}

We record explicitly why the remainder of UT is valid for these
arbitrary companions. Define the specific sorted polynomial
\(B_b^\pi\) and its polarization \(\mathcal T_b^\pi\) by UT2.
The invariant Courant argument in UT2--UT3 was already proved
for arbitrary \(\mathcal D(Y_U)\)-inputs: its auxiliary invariant
finite PL expressions lie in \(\mathcal D\) by lattice closure,
and FS descends them. Its error terms are literal full-carrier
zero-product generators, even after multiplication by arbitrary
downstream tails.

Fix \(F,L\), and put \(C=F\cap L\). If \(C\) is empty,
all relevant KC costs are zero. Otherwise fix
\(v=p-r\in\mathcal D(Y_U)\) with one global original psh
difference decomposition. Let \(\Delta_{I,k}\) be the protected
prefix minors of \(e\), and set
\[
 \begin{aligned}
 D_C&=\max_{k,I}\|\log|\Delta_{I,k}|\|_{q^{-1}(C)},\\
 M_C&=d\bigl(\|p\|_{\pi^{-1}(C)}
                    +\|r\|_{\pi^{-1}(C)}\bigr)+D_C.
 \end{aligned}                                               \tag{LC24.8}
\]
The lifts are complete compact inverse images. For \(N\ge1\),
the invariant maxima
\[
 \max_{|I|=k}\left(
     \sum_{i\in I}p_i+\sum_{i\notin I}r_i
                      +N^{-1}\log|\Delta_{I,k}|\right)
                                                               \tag{LC24.9}
\]
are finite continuous psh functions on the ordering cover and
descend psh by FS. At \(k=0\) use \(\operatorname{Tr}_\pi r\)
and the empty minor of logarithm zero. Their differences give
\(N^{-1}(L_k^e(Nv)-L_{k-1}^e(Nv))\).
Omitting the minor term gives the corresponding order statistic
\(a_k(v)\). Their uniform difference on \(C\) is at most
\(2D_C/N\), and both chosen psh-difference decompositions have
total centered oscillation cost at most \(2M_C\).
These facts use only \(v=p-r\), not a finite PL presentation.

If \(R=\prod_{j=1}^{\ell}[z_j]\), choose each fixed
\(z_j=p_j-r_j\), and put
\(J_j=\operatorname{osc}_C p_j+\operatorname{osc}_C r_j\).
The exact KC telescope UT12--UT14 gives
\[
 q^{U,b+\ell}_{F,L}
 \left(\bigl(N^{-b}C_b^e(Nv)-B_b^\pi(v)\bigr)R\right)
 \le
 \frac{2dbD_C}{N}(2M_C)^{b-1}\prod_{j=1}^{\ell} J_j.              \tag{LC24.10}
\]
Here \(e\) in \(C_b^e\) denotes the fixed frame and
\(\ell=\deg R\).
Empty carrier stages have zero cost; zero oscillation budgets
and \(b=1\) have the same conventions as UT14.

Put \(h=\log|g|\). The conditional UT7 just proved can be
iterated at the weights \(Nv+jh\) for every integer \(j\):
all these weights belong to \(\mathcal D(Y_U)\).
For each fixed integer \(M\), forward or reverse telescoping
gives
\[
 q^{U,b+\deg R}_{F,L}
 \left(
  \bigl(C_b^e(N(v+Mh))-C_b^e(Nv)\bigr)R
 \right)=0 .
\]
Apply (LC24.10) to the two fixed functions \(v+Mh,v\),
and let \(N\to\infty\). This yields
\[
 q^{U,b+\deg R}_{F,L}
        \left((B_b^\pi(v+Mh)-B_b^\pi(v))R\right)=0
                    \quad(M\in\mathbf Z).                 \tag{LC24.11}
\]
The same initial \(U\), frame and finite gallery are used for
all \(N,M,v\). No limit in the analytic relation ideal is asserted;
the conclusion follows from the displayed numerical bound in
the unchanged seminorm.

The Vandermonde inversion at \(M=1,\ldots,b\), using UT3,
gives
\(q(\mathcal T_b^\pi(h,v,\ldots,v)R)=0\).
Finite polarization in \(v\) gives the conclusion for unrelated
\(v_2,\ldots,v_b\in\mathcal D(Y_U)\). The zero-product
certificate errors remain legal after \(R\).
For \(b=1\), the trace is the global harmonic function
\(\log|\operatorname{Norm}_\pi g|\), so the result is direct.

\begin{proposition}[Conditional arbitrary-companion unit transfer]
For the fixed original finite \'{e}tale cover, point \(x\) and
global unit \(g\), choose the GP representative \(U\) before all
weights, degrees, carriers, stages, polars and downstream tails.
If (LC24.2) holds on this same \(U\), then for every \(b\ge1\),
every \(v_2,\ldots,v_b\in\mathcal D(Y_U)\), every closed
\(F\subset U\), every compact \(L\Subset U\), and every finite
formal downstream monomial \(R\) in \(\mathcal D(U)\),
\[
 q^{U,b+\deg R}_{F,L}
       \bigl(\mathcal T_b^\pi(\log|g|,v_2,\ldots,v_b)R\bigr)=0.
                                                               \tag{LC24.12}
\]
The degree-one case does not require the assumption.
\end{proposition}

The proof above is the complete deduction. When \(F\) is an
actual absolute cap carrier, it is kept unchanged. KC6 gives
the action consequence in all permitted degrees for every
original symmetric middle-positive separately closed background,
all original smooth complementary tests, and every actual
refinement. This is only the forward consequence of the
conditional certificate; no new positive-current pullback is used.

\subsection{Remaining gate and relation to the branch obstruction}

BC shows that the global-harmonic-comparator replacement for
(LC24.2) cannot hold uniformly for all upstairs weights, even
for one fixed quadratic cover and one fixed unit. The present
proof asks instead for local annihilation of the already existing
original Radon measure. The strict-defect function need not
extend harmonically to \(U\). Thus BC does not refute the
antecedent of this proposition.

Conversely, this proposition does not prove that antecedent.
It does not transfer an arbitrary upstairs function whose trace
is merely locally harmonic along the whole lifted carrier:
such a function need not be one global unit logarithm.
Nor does it transfer every allowed ambient generator in all its
original local-presentation quantifiers. Therefore full ST19,
finite-cover reverse carrier-CLN comparison, LG7, KC12 and CTO
remain \texttt{OPEN}. The result is a conditional implication
between explicit existing gaps, not a new abstract definition
or a general annihilation theorem.

\paragraph{Construction and review roles.}
The root supplied the conditional local-measure replacement for
the BC-false global-comparator route. The polar-duality researcher
checked and wrote its complete strict-defect descent, all-tail
exchange and quantified UT deduction. Neither is an independent
reviewer of this source. Their earlier authorship of GP, CW,
UT, IR/IL/NC is input provenance only. Any author-separated
review of this new deduction is recorded against exact bytes.

